% =============================================================================
% Section IV: Decoherence in Power-Law FLRW Spacetimes
% paper/powerlaw.tex
% =============================================================================

\section{Decoherence in Power-Law FLRW Spacetimes}
\label{sec:powerlaw}

We now evaluate the decoherence exponent $\Gamma$ — equivalently the decoherence
number $\Nnum = \Gamma/2$ — for the power-law scale factor
$\Scfac(\tau) = C_{p}(\tau - \tau_{*})^{p}$ with $p > 1$, for which a future
causal horizon is present (\cref{sec:flrw-powerlaw}).  The superposition is
switched on adiabatically at proper time $\tau_{\rm on}$ and held for proper
duration $T$ with fixed separation $d(\tau) = d_{0}$.

\subsection{Reduction to a one-dimensional integral}
\label{sec:powerlaw-setup}

Starting from the dipole decoherence formula~\cref{eq:Gamma-dipole}, we change
the integration variable from proper time $\tau$ to conformal time $t$ using
$d\tau = \Scfac(t)\,dt$.  Writing the separation vector as $s^{a} = d_{0}\,\hat{s}^{a}$
(with $\hat{s}^{a}$ a unit comoving spatial vector), the proper-time integral becomes
%
\begin{align}
  \Gamma
  &=
  q^{2}
  \int_{t_{\rm on}}^{t_{\rm off}}\!\!\!dt_{1}
  \int_{t_{\rm on}}^{t_{\rm off}}\!\!\!dt_{2}\;
  d_{0}^{2}\,\Scfac(t_{1})\,\Scfac(t_{2})\;
  \hat{s}^{a}\hat{s}^{b}
  \,\pd_{a}\pd_{b}\,
  \Wight\!\lb(t_{1}, \mb{X};\, t_{2}, \mb{X}\rb) \,,
  \label{eq:Gamma-conformal}
\end{align}
%
where $\pd_{a}$ now denotes a derivative with respect to comoving spatial
coordinates (no factors of $\Scfac$ in the comoving metric).  Using the
conformal-vacuum Wightman function~\cref{eq:Wight-FLRW} at coincident spatial
position $\mb{x}_{1} = \mb{x}_{2} = \mb{X}$,
%
\begin{equation}
  \Wight(t_1, \mb{X};\, t_2, \mb{X})
  =
  \frac{-1}{4\pi^{2}\,\Scfac(t_{1})\,\Scfac(t_{2})
    \,\tilde{\sigma}_{\varepsilon}(t_{1}, t_{2})} \,,
  \label{eq:Wight-coincident}
\end{equation}
%
with $\tilde{\sigma}_{\varepsilon}(t_{1},t_{2}) = -(t_{1}-t_{2}-i\varepsilon)^{2}$
at $\mb{x}_{1} = \mb{x}_{2}$.  Taking two spatial derivatives along the separation
direction and evaluating at coincident spatial position gives (see \cref{app:integral}
for the explicit computation)
%
\begin{equation}
  \hat{s}^{a}\hat{s}^{b}\,\pd_{a}\pd_{b}\,
  \Wight\!\lb(t_{1}, \mb{X};\, t_{2}, \mb{X}\rb)
  =
  \frac{-1}{4\pi^{2}\,\Scfac(t_{1})\,\Scfac(t_{2})}
  \cdot
  \frac{1}{3}\,\partial_{t_{1}}^{2}
  \lb[\tilde{\sigma}_{\varepsilon}(t_{1},t_{2})^{-1}\rb] \,.
  \label{eq:dipole-kernel}
\end{equation}
%
The factor of $1/3$ arises from averaging the squared spatial projection over
the two transverse directions, valid for an isotropic dipole orientation (see
\cref{app:dipole}).  Substituting~\cref{eq:dipole-kernel}
into~\cref{eq:Gamma-conformal} and integrating the resulting expression, the
two conformal factors $\Scfac(t_{1})\Scfac(t_{2})$ from
$d\tau_{1}d\tau_{2} = \Scfac(t_{1})\Scfac(t_{2})\,dt_{1}dt_{2}$ precisely cancel
those in the denominator of $\Wight$ at coincident position,
yielding the central integral of this paper:
%
\begin{align}
  \Gamma
  &=
  \frac{q^{2}d_{0}^{2}}{12\pi^{2}}
  \,\mathrm{Re}
  \int_{t_{\rm on}}^{t_{\rm off}}\!\!\! dt_{1}
  \int_{t_{\rm on}}^{t_{\rm off}}\!\!\! dt_{2}\;
  \frac{1}{\lb(t_{1} - t_{2} - i\varepsilon\rb)^{4}} \,,
  \label{eq:decoherence-integral}
\end{align}
%
where $\varepsilon \to 0^{+}$ and the real part selects the physical (symmetric)
contribution.  The conformal time limits are
%
\begin{equation}
  t_{\rm on}
  =
  \int_{\tau_{*}}^{\tau_{\rm on}} \frac{d\tau}{\Scfac(\tau)}
  =
  \frac{u_{\rm on}^{1-p}}{C_{p}(p-1)} \,,
  \qquad
  t_{\rm off} = t_{\rm on} + \Delta t \,,
  \label{eq:conformal-limits}
\end{equation}
%
where $u_{\rm on} = \tau_{\rm on} - \tau_{*}$ and $\Delta t = \int_{\tau_{\rm on}}^{\tau_{\rm on}+T} d\tau/\Scfac(\tau)$ is the conformal-time duration of the superposition.  Because $p > 1$, the conformal time $t_{\rm off}$ remains bounded below the finite horizon value $t_{\infty} = t_{\rm on} + u_{\rm on}^{1-p}/[C_{p}(p-1)]$ for any finite $T$.

The double integral~\cref{eq:decoherence-integral} depends on the background
geometry only through the conformal time limits $[t_{\rm on}, t_{\rm off}]$.  No
factor of $\Scfac$ appears in the integrand: the conformal invariance of the
scalar field is fully manifest.

\subsection{The decoherence exponent: closed-form result}
\label{sec:powerlaw-result}

The integral in~\cref{eq:decoherence-integral} is evaluated in
\cref{app:integral}.  The result depends on a single dimensionless ratio
$\widetilde{T} = T/T_{0}$, where
%
\begin{equation}
  T_{0}
  \equiv
  \frac{p-1}{p}\,
  \frac{u_{\rm on}}{C_{p}\,u_{\rm on}^{p}}
  =
  \frac{u_{\rm on}^{1-p}}{C_{p}(p-1)}\cdot\frac{(p-1)^{2}}{p}
  \,,
  \label{eq:T0-def}
\end{equation}
%
is a characteristic proper-time scale set by the scale factor at the onset of
the superposition.  Physically, $T_{0}$ is of order the proper time remaining
until the future causal horizon is reached from $\tau_{\rm on}$: a superposition
of duration $T \sim T_{0}$ has the source radiating for a conformal-time interval
comparable to the remaining conformal-time range.  Explicitly, in terms of
$u_{\rm on}$ and the scale-factor parameters,
%
\begin{equation}
  T_{0} = \frac{(p-1)^{2}}{p}\,
  \frac{u_{\rm on}^{1-p}}{C_{p}(p-1)}
  = \frac{(p-1)}{p}\,
  \frac{1}{C_{p}}\, u_{\rm on}^{1-p} \,.
  \label{eq:T0-explicit}
\end{equation}
%

The closed-form result for the decoherence number $\Nnum = \Gamma/2$ is
%
\begin{equation}
  \boxed{
  \Nnum(T)
  =
  \frac{q^{2}d_{0}^{2}\,p^{3}}{48\pi^{3}}
  \cdot
  \frac{2\,T_{0}^{2} + \widetilde{T}\,(2 + \widetilde{T})}{T_{0}^{2}\,(1 + \widetilde{T})^{2}}
  \cdot
  \ln(1 + \widetilde{T})
  }
  \label{eq:N-exact}
\end{equation}
%
where $\widetilde{T} = T/T_{0}$.  Equation~\eqref{eq:N-exact} is exact for any
$T > 0$ and any $p > 1$.  Several structural features are worth noting
immediately.  First, $\Nnum(0) = 0$ as required: no decoherence before the
superposition is created.  Second, $\Nnum(T) > 0$ for all $T > 0$, confirming
that decoherence begins immediately once the superposition is switched on.
Third, the formula contains a single overall scale $q^{2}d_{0}^{2}p^{3}$ and a
single dimensionless argument $\widetilde{T}$, so the entire $T$- and
$p$-dependence is captured by two parameters.

\subsection{Analysis of the result}
\label{sec:powerlaw-analysis}

\paragraph{Large-$T$ behavior.}
For $T \gg T_{0}$ (i.e., $\widetilde{T} \gg 1$), the rational prefactor
$(2T_{0}^{2} + \widetilde{T}(2+\widetilde{T}))/[T_{0}^{2}(1+\widetilde{T})^{2}]
\to T_{0}^{-2}$, and therefore
%
\begin{equation}
  \Nnum(T)
  \;\xrightarrow{T \gg T_{0}}\;
  \frac{q^{2}d_{0}^{2}\,p^{3}}{48\pi^{3}\,T_{0}^{2}}
  \,\ln\!\lb(\frac{T}{T_{0}}\rb) \,,
  \label{eq:N-large-T}
\end{equation}
%
up to corrections of order $T_{0}/T$.  The decoherence grows logarithmically
without bound: $\Nnum \to \infty$ as $T \to \infty$.  For any nonzero charge and
separation, the superposition is therefore \emph{inevitably} fully decohered
at sufficiently late times.  The logarithmic growth is slower than the linear
growth $\Nnum \sim T$ found in de Sitter spacetime~\cite{Danielson:2022tdw}, but
both are unbounded, and in both cases the divergence is a direct consequence of the
causal horizon: soft radiation irreversibly crosses the horizon, carrying
which-path information to permanently inaccessible regions.  We return to this
comparison in \cref{sec:desitter}.

\paragraph{Early-time behavior.}
For $T \ll T_{0}$ (i.e., $\widetilde{T} \ll 1$), expanding to leading non-trivial
order,
%
\begin{equation}
  \Nnum(T)
  \;\xrightarrow{T \ll T_{0}}\;
  \frac{q^{2}d_{0}^{2}\,p^{3}}{48\pi^{3}}
  \cdot
  \frac{T^{3}}{3\,T_{0}^{4}} \,,
  \label{eq:N-small-T}
\end{equation}
%
so $\Nnum \propto T^{3}$ at early times.  This power-law growth reflects the
short-distance behavior of the Wightman function before the superposition has
time to probe the causal structure of the horizon.  Physically, at early times the
radiation has not yet had time to reach the horizon; the system behaves as if it
were in a slowly evolving Minkowski background, and the decoherence accumulates
as a local process.  The crossover from cubic to logarithmic growth occurs
at $T \sim T_{0}$, when the radiation wavelengths become comparable to the
remaining conformal-time distance to the horizon.

\paragraph{Dependence on $p$.}
The overall prefactor $p^{3}$ in~\cref{eq:N-exact} shows that the decoherence
rate increases rapidly with the expansion rate exponent.  A spacetime with
$p = 5$ (strongly accelerated expansion) decoheres $125$ times faster than one
with $p = 1$ at the same values of $q$, $d_{0}$, and $T_{0}$.  In addition,
$T_{0}$ itself depends on $p$ through~\cref{eq:T0-explicit}: faster expansion
($p$ larger) moves the horizon closer in proper time, so $T_{0}$ is smaller.
Both effects conspire to make the onset of strong decoherence earlier as $p$
increases.

The numerical behavior of $\Nnum(T)/\Nnum(T_{0})$ as a function of $\widetilde{T}$
for representative values $p = 2, 3, 5$ is shown in \cref{fig:decoherence},
together with the de Sitter linear growth for comparison.  All three power-law
curves exhibit the same qualitative shape — cubic onset transitioning to
logarithmic growth — but the normalized amplitude $\Nnum(T_{0})/\Nnum^{(p=2)}(T_{0})$
differs between curves due to the $p^{3}$ factor.

\subsection{The no-horizon case: $p = 1$}
\label{sec:powerlaw-nohorizon}

The qualitative distinction between the $p > 1$ (horizon) and $p = 1$ (no horizon)
cases is the central diagnostic of the causal mechanism.  For $p = 1$, the
conformal time is
%
\begin{equation}
  t(\tau) = \frac{1}{C_{1}}\ln\!\lb(\frac{\tau - \tau_{*}}{u_{\rm on}}\rb) \,,
  \label{eq:t-pone}
\end{equation}
%
which runs from $0$ to $+\infty$ as $\tau$ runs from $\tau_{\rm on}$ to
$\infty$.  The conformal time range is therefore infinite: no future causal
horizon exists (\cref{sec:flrw-powerlaw}).

In this case the upper limit $t_{\rm off} = t_{\rm on} + \int_{\tau_{\rm on}}^{\tau_{\rm on}+T} d\tau/\Scfac(\tau)$
grows without bound as $T \to \infty$, and the double integral
in~\cref{eq:decoherence-integral} must be re-evaluated.  In the $p \to 1$
limit, the rational-times-logarithm structure of~\cref{eq:N-exact} changes
character: the conformal-time duration $\Delta t \to \infty$ does not introduce
an IR divergence in $\Nnum$ because the integrand of~\cref{eq:decoherence-integral}
is a power-law $\sim(t_{1}-t_{2})^{-4}$, which is \emph{integrable} for the
compact differences that arise.  One can show (see \cref{app:integral}) that
the $T \to \infty$ limit of the double integral is finite:
%
\begin{equation}
  \Nnum(T \to \infty)\big|_{p=1}
  =
  \frac{q^{2}d_{0}^{2}}{16\pi^{2}\,u_{\rm on}^{2}} \,,
  \label{eq:N-nohorizon}
\end{equation}
%
where $u_{\rm on} = \tau_{\rm on} - \tau_{*}$ is the proper time elapsed since
the Big Bang at the start of the superposition.  The decoherence \emph{saturates}
to the finite value~\cref{eq:N-nohorizon} as $T \to \infty$: the off-diagonal
element $|\langle \gamma_{1}|\rho_{\rm red}|\gamma_{2}\rangle|$ approaches the
constant $\exp(-\Nnum_{\infty})$ and does not vanish.

The physical interpretation is transparent.  For $p = 1$ there is no future
causal horizon, so the soft radiation emitted by the superposed charge propagates
to spatial infinity but is not permanently hidden from any comoving observer.
In principle, an experimenter with access to the full future null infinity
$\ms{I}^{+}$ could measure the outgoing radiation and recover the which-path
information — the field state remains pure when all of $\ms{I}^{+}$ is accessible.
The saturation of $\Nnum$ reflects exactly this: the information is \emph{displaced}
(encoded in the radiation field) but not \emph{destroyed}.

Contrasting with the $p > 1$ case: the causal horizon acts as a one-way membrane
that sequesters a portion of $\ms{I}^{+}$ from the comoving observer, permanently
removing its contribution from the accessible Hilbert space.  That irreversible
loss is what converts a saturating $\Nnum$ into a diverging one, and it is the
causal horizon — not the dynamics of the scale factor per se — that is the
essential physical ingredient.  The $p = 1$ case makes this precise: the scale
factor $\Scfac(\tau) = C_{1}(\tau - \tau_{*})$ also describes a non-trivial
expanding universe (linear expansion), yet without a horizon, complete decoherence
is never achieved.

%% CITATIONS NEEDED (for gpd-bibliographer):
%% - Danielson:2022tdw  -- DSW (2023) PRD 108:025007, arXiv:2301.00026 (de Sitter N~H^3 T result)
%% - Danielson:2024yge  -- DSW (2025a) PRD 111:025014, arXiv:2407.02567 (local framework)
